\chapter{Writing Clear Scientific Prose}
\label{chap:text}
\index{typography}\index{emphasis}\index{quotation marks}\index{footnotes}\index{special characters}

A document can compile perfectly and still be hard to read. Scientific prose
needs visible paragraphs, purposeful emphasis, accurate symbols, and source
that another author can revise without decoding improvised formatting. In this
chapter we give ordinary text the care that mathematical notation will receive
later.

\begin{learningoutcomes}
  \item Create paragraphs and distinguish them from line breaks.
  \item Apply emphasis, bold, and italic text for stated purposes.
  \item Type quotation marks, dashes, ellipses, and reserved characters.
  \item Insert non-breaking spaces, footnotes, and web addresses.
  \item Format a biological species name conventionally.
  \item Explain a safe beginner approach to Unicode source.
\end{learningoutcomes}

\section{Paragraphs carry ideas}

A paragraph is a unit of thought, not merely a block that fits a page. In
source, leave one blank line when the argument moves to a new paragraph.
\LaTeX{} decides the printed line endings.

\begin{latexcode}{Complete example: two paragraphs}
\documentclass{article}
\begin{document}
Method A records one value for each specimen. Method B records
the same property independently.

We compare paired values because both methods examine the same
specimens.
\end{document}
\end{latexcode}

The PDF contains two paragraphs. The single source line ending inside each
paragraph becomes an ordinary space. This separation lets source lines remain
short enough to review.

The command \code{\\} requests a line break without starting a new paragraph.
It is useful in structures such as an address or a short title block, but it is
usually wrong between prose paragraphs. A blank line expresses the meaning
more accurately.

\begin{commonerror}
\textbf{Symptom:} prose has a ragged sequence of short printed lines.

\textbf{Cause:} every source line ends with \code{\\}.

\textbf{Correction:} remove the forced breaks and let \LaTeX{} compose the
paragraph. Keep a blank source line only where a new paragraph begins.
\end{commonerror}

\section{Emphasis and type styles}

Use \cmd{emph}\required{text} when the meaning is emphasis. In ordinary text,
\LaTeX{} generally renders it in italics; inside existing emphasis it can switch
back, which preserves the distinction.

\begin{latexcode}{Complete example: emphasis and styles}
\documentclass{article}
\begin{document}
The values are \emph{simulated}, not observations.

\textbf{Warning:} do not interpret them as clinical evidence.

The species name \textit{Drosophila melanogaster} is italic.
\end{document}
\end{latexcode}

\cmd{textbf} requests bold text; use it for a genuine visual role such as a
short warning label, not for every important sentence. \cmd{textit} requests
italic shape without claiming emphasis, which suits terms conventionally set
in italics. Semantic structure remains preferable: a section title should be
made with \cmd{section}, not with a large bold line.

\begin{scienceinpractice}
Biological genus and species names are conventionally italic, while the
authority is not: \textit{Panthera leo} (Linnaeus, 1758). After the first full
mention, a field may permit an unambiguous abbreviation such as
\textit{P.~leo}. Follow the relevant disciplinary style and avoid automatic
abbreviation when two genera could share the initial.
\end{scienceinpractice}

\section{Quotation marks, hyphens, and dashes}

In pdf\LaTeX{} source, type an opening double quotation as two grave accents and a
closing one as two apostrophes. Single quotations use one of each:

\begin{latexcode}{Complete example: quotations and dashes}
\documentclass{article}
\begin{document}
The manual calls this a ``minimal working example''.
The term `minimal' does not mean incomplete.

A high-resolution image uses a hyphen. Pages 12--18 use an
en dash. A sudden interruption---used sparingly---uses an em dash.
\end{document}
\end{latexcode}

One hyphen joins a compound; two hyphens in source produce an en dash for a
range or relationship; three produce an em dash for a break in thought.
British editorial styles differ over spaces around an em dash. This book uses
none and applies the choice consistently.

Use \cmd{ldots} for a deliberate ellipsis in text: \ldots. Three full stops
may have uneven spacing. An ellipsis signals omitted or trailing material; it
should not become a decorative ending for uncertain prose.

\begin{pausepredict}
Predict the output of \code{10-12}, \code{10--12}, and \code{10---12} before
compiling them. Which source represents a numerical range?
\end{pausepredict}

\section{Spaces that should stay together}

A tilde in prose creates a \term{non-breaking space}. It prevents a line break
between related items while printing an ordinary-looking space:

\begin{latexcode}{Useful non-breaking spaces}
Figure~3
Dr~Rao
Section~2.1
\end{latexcode}

These fragments are not a complete document, but each shows a relationship
that should not split across lines. Later, automatic cross-references will
replace typed numbers such as \code{Figure~3}. Do not insert tildes between
every pair of words; excessive restrictions can produce poor spacing.

\section{Footnotes and web addresses}

A footnote command contains the note where its marker belongs:

\begin{latexcode}{Complete example: a footnote and URL}
\documentclass{article}
\usepackage{hyperref}
\begin{document}
The data are simulated.\footnote{The file is supplied only for
typesetting practice.}

Documentation is available at
\url{https://www.latex-project.org/help/documentation/}.
\end{document}
\end{latexcode}

\cmd{footnote} prints a marker and places the note at the bottom of the page.
Use footnotes for short supplementary remarks, not for essential methods or a
hidden chain of citations. The \pkg{hyperref} package provides \cmd{url},
which allows a web address to break more safely and makes it a link in the PDF.

Long or unstable URLs are rarely good citation substitutes. A bibliography
record can carry the author, title, date, access date, and persistent
identifier. Chapter~\ref{chap:citations} develops that workflow.

\section{Unicode without mystery}

Modern \LaTeX{} source is normally saved as UTF-8, an encoding capable of
representing many writing systems. A current engine can read common accented
Latin characters directly, but font and engine support still matter. A file
copied from a website may also contain invisible spaces or typographic quotes
that look similar to ASCII source syntax.

For beginner examples we keep command syntax and quotation marks in ASCII. An
accent command such as \code{Jos\textbackslash\'e} is portable in a small
pdf\LaTeX{} example. Later, a project requiring extensive multilingual text may
choose Lua\LaTeX{} or Xe\LaTeX{} with suitable fonts after checking the publisher's
workflow.

\begin{goodpractice}
Do not repair a mysterious character by repeatedly changing the compiler.
Identify the character, confirm that the file is UTF-8, and test a minimal
example. A journal template may require a specific engine.
\end{goodpractice}

\chapterproject

\begin{miniproject}
Improve the prose in the running article. Add two background paragraphs,
keeping the existing simulated-data statement. Use \cmd{emph} once to stress
that the values are simulated, type \code{Method A\textbackslash\&B} only if an
ampersand is genuinely part of a name, and add a short footnote explaining
that Method A and Method B are neutral teaching labels. Compile and check that
no forced line breaks interrupt the prose.
\end{miniproject}

\chaptersummary

Blank source lines make paragraphs; ordinary source line endings do not force
printed breaks. \cmd{emph} expresses emphasis, while \cmd{textbf} and
\cmd{textit} request particular shapes. ASCII quotation conventions, source
hyphens, \cmd{ldots}, non-breaking spaces, escaped characters, footnotes, and
\cmd{url} give prose precise typography. UTF-8 is the normal source encoding,
but unusual characters still need compatible fonts and engines.

\chaptercheck
\begin{chapterchecklist}
  \item I use blank lines, not repeated \code{\\}, for prose paragraphs.
  \item I can explain why a word is emphasised, bold, or italic.
  \item I can type opening and closing quotation marks in ASCII source.
  \item I choose the right source form for a hyphen, range, or em dash.
  \item I can add a footnote and a safe URL.
  \item I have checked the running article's paragraph structure.
\end{chapterchecklist}

\chapterexercisestitle
\begin{chapterexercises}
  \item Explain the difference between a source line ending and a blank line.
  \item Correct a paragraph that ends every source line with \code{\\}.
  \item Write source for the quotation: The result was ``unexpected''.
  \item Choose a hyphen or dash for each: pages 4 to 9; evidence based; an
        interruption in a sentence.
  \item Give one suitable and one unsuitable use of a footnote.
  \item Format \textit{Escherichia coli} and its later abbreviation.
  \item Debug source in which \code{50\%} accidentally comments out the rest
        of the line.
  \item Challenge: edit a paragraph from your field using only semantic
        emphasis and explain every typographic choice.
\end{chapterexercises}
